\documentclass[UTF8,fontset=windows,11pt]{ctexart}
\usepackage[a4paper,margin=2.1cm]{geometry}
\usepackage{amsmath,amssymb}
\usepackage{booktabs}
\usepackage{tabularx}
\usepackage{longtable}
\usepackage{array}
\usepackage{ragged2e}
\usepackage{enumitem}
\usepackage{xcolor}
\usepackage{tikz}
\usetikzlibrary{positioning,arrows.meta,calc}
\usepackage{titlesec}
\usepackage{titling}
\usepackage[colorlinks=true,linkcolor=NavyBlue,urlcolor=NavyBlue,citecolor=NavyBlue]{hyperref}

\definecolor{NavyBlue}{RGB}{30,70,140}
\definecolor{AcademicRed}{RGB}{170,40,40}
\definecolor{boxgray}{RGB}{245,245,247}

\newcolumntype{Y}{>{\RaggedRight\arraybackslash}X}
\newcommand{\key}[1]{\textbf{\textcolor{NavyBlue}{#1}}}

\xeCJKsetcharclass{"2460}{"24FF}{1}
\xeCJKsetcharclass{"2605}{"2606}{1}

\setlist[itemize]{leftmargin=1.4em,topsep=2pt,itemsep=1pt}
\setlist[enumerate]{leftmargin=1.6em,topsep=2pt,itemsep=1pt}

\titleformat{\section}{\Large\bfseries\color{NavyBlue}}{\thesection}{0.6em}{}
\titleformat{\subsection}{\large\bfseries\color{NavyBlue!85!black}}{\thesubsection}{0.6em}{}

\tikzset{
  band/.style={rounded corners=3pt, draw=black!35, line width=0.5pt, text width=15.3cm, inner sep=6pt, align=left, font=\small},
  flow/.style={-{Stealth[length=2.4mm]}, line width=0.7pt, black!55},
}

\title{\vspace{-1.2cm}\textbf{SYNAPSE · 协作即压缩}\\[2mm]
\large 带增长记忆的多智能体语义 Wyner--Ziv 通信 · 设计与执行方案}
\author{}
\date{内部设计稿 · 2026-06-18}

\begin{document}
\maketitle
\vspace{-1.0cm}

\begin{quote}\small\color{black!80}
\key{协作即压缩（Coordination as Compression）}：把"两个智能体达成可继续任务的对齐"建模为一个带\textbf{可学习、随经验增长的解码端边信息}的 Wyner--Ziv 编码问题；由此得到一条可测、可证伪的\textbf{协作速率收缩律}——并在分布漂移时回弹，使速率本身成为一个免费的传感器。

本方案历经多轮自审计与红队迭代收敛定稿：一份既能拿下竞赛国一、又能投 CCF-A 主会的完整设计、实验与执行方案。

一句话定位：\textbf{别去抢"第一个把率失真用于多智能体"（已有 arXiv:2604.09521 占了静态能力角度）；去抢谁都还没占的"第一个证明协作速率有一条会收缩、会回弹、还能当漂移传感器的时间律，并在真实 openEuler 系统里跑出来"。}
\end{quote}

\section*{〇、给三类读者的三句话}
\addcontentsline{toc}{section}{〇、给三类读者的三句话}
\begin{itemize}
\item \key{给评委}：我们做的不是"又一个省 token 的小聪明"。我们让一群智能体\textbf{越合作越少说废话}——每个新任务需要在链路上交换的字节，沿一条曲线下降、逼近物理下限；一旦任务类型突变，字节又会\textbf{当场回弹}。这条曲线我们用 eBPF \textbf{实测}给你看，而且有损压缩\textbf{保证不出错}（带校验和）。
\item \key{给审稿人}：我们把多智能体对齐形式化为\textbf{带增长边信息的语义 Wyner--Ziv}，给出可达方案与失配惩罚 $\Delta$，并提出一个\textbf{非显然、可证伪的预测}——收缩曲线会停在一个可独立测量的信息地板上。漂移回弹给出一个\textbf{模型无关的分布漂移传感器}。
\item \key{给自己（工程纪律）}：一套系统，两层产出。\textbf{L0 全确定性、现场能稳跑}（拿竞赛分）；\textbf{L1 学习层离线出曲线}（成论文）。同源，不重复造轮子。
\end{itemize}

\section{研究背景：多智能体协作，贵在"反复说"}

\subsection{问题}
LLM 多智能体系统（MAS）正在从"单体大模型"走向"分工协作"。但协作的代价被长期低估：\textbf{每多一个 agent、每多一轮交互，就要把上下文反复序列化成文本、再解析回来}。在长上下文、多角色、连续任务的真实工作流里，\textbf{通信开销吃掉的算力与时延，常常比任务本身还多}。学界已有三条线在缓解它，但各有缺口：

\begin{table}[h]\footnotesize\centering
\renewcommand{\arraystretch}{1.25}
\begin{tabularx}{\linewidth}{@{}p{2.1cm}p{3.2cm}Y Y@{}}
\toprule
\textbf{路线} & \textbf{代表} & \textbf{在省什么} & \textbf{缺口} \\
\midrule
潜空间通信 & CIPHER (2310.06272)、LatentMAS+OBF、KVComm (2510.03346) & 发潜表示替代文本 & 不可审计、要求同构、\textbf{会静默损坏}、且潜载荷按字节常\textbf{更大} \\
缓存 & Agentic Plan Caching (2506.14852)、语义缓存 & 重复再生成的算力（compute） & 只对\textbf{重复}任务有效；命中率低则失效；省的是"重算"，不是"重说" \\
共享记忆 & Collaborative Memory (2505.18279) & 把记忆当外挂模块 & 与通信各用一套机制；读写多为规则式、静态 \\
\bottomrule
\end{tabularx}
\end{table}

\subsection{两个被忽略的轴}
现有"通信效率"研究几乎都在问\textbf{空间轴}的问题：开销如何随 agent 数 $N$ 增长（拓扑剪枝、AgentPrune 2410.02506）。\textbf{没有人系统地刻画时间轴}：当一支团队\textbf{跨任务累积共享经验}后，单位任务的协作开销如何\textbf{随经验衰减}？它有没有下限？会不会回弹？

而最近一篇理论工作 \key{arXiv:2604.09521《Semantic Rate-Distortion for Bounded Multi-Agent Communication》} 已经把率失真/Wyner--Ziv 搬进多智能体通信——但它处理的是\textbf{静态的能力异质}（不同算力的 agent 需要不同的语义字母表），\textbf{纯理论、无系统、无时间动力学}。

\begin{quote}\small
\textbf{这恰好把我们的位置照亮了}：率失真/WZ 作为"透镜"已被学界认可（去风险），但\textbf{"边信息随经验增长 $\to$ 速率随时间收缩 $\to$ 漂移回弹"这一整块动力学，加上真实系统的字节级实证，仍是无人区。}这就是 SYNAPSE 要占的山头。
\end{quote}

\subsection{为什么现在}
2025--2026 是 agent 通信的爆发窗口（C2C 2510.19995 把通信成本显式建模、Scaling Teams or Scaling Time 2604.03295 研究终身学习 MAS 的性能-时间权衡）。但它们都没回答"\textbf{协作通信速率随时间怎么走}"。窗口开着，且正在收窄——这是一个值得现在做的题。

\section{核心思想：协作即压缩}

\subsection{直觉：老同事为什么话越来越少}
新搭档时，每件事都要从头交代："那个客户叫什么、上次为什么砍预算、模板在哪。"合作久了，你只要说一句"老规矩，把 Q3 那版改一下发客户"，对方全懂。

\textbf{话变少，不是因为你把上次的话录下来重放（那是缓存）}，而是因为\textbf{对方脑子里长出了一个能预测你意图的模型}。你只需要说\textbf{对方猜不到的那一点点新信息}，其余他自己补全。SYNAPSE 把这件事搬到 LLM 智能体之间，并且\textbf{量化}它：

\begin{enumerate}
\item \key{只发"惊讶"}——发送方不发完整内容，只发\textbf{接收方用自己的记忆预测不出来的残差}。对方能猜到的，一个字都不发。
\item \key{边信息会长大}——"对方能预测的部分"随\textbf{共享记忆跨任务巩固}而越来越大，于是\textbf{每个新任务要传的字节越来越少}。这就是\textbf{收缩律}。
\item \key{有一个不可逾越的地板}——话不会变成零。每个新任务总有\textbf{真正新的、谁都猜不到的信息量} $H(Y\mid B_\infty)$。这是物理下限。
\item \key{换了赛道，话又变多}——任务类型突变（分布漂移），老经验失效，\textbf{字节数当场回弹}；而且回弹发生在\textbf{任务质量掉下去之前}。所以"话突然变多"本身就是一个\textbf{免费的、提前的"风向变了"警报}。
\item \key{有损但保证不出错}——残差是有损压缩，万一对方猜错？\textbf{每条消息附几十字节的哈希校验}，对方重构完一对，不符就回退取全量。\textbf{省了带宽，端到端绝不静默出错}。
\end{enumerate}

\subsection{形式化：Wyner--Ziv 透镜 + 收缩律}
\textbf{设定}。$N$ 个 agent 协作完成任务流 $\{T_1,T_2,\dots\}\sim\mathcal D$。共享记忆 $\mathcal M$ 是内容寻址的类型化状态库（\textbf{既是通信媒介，又是持久记忆}——这是 v1"突触"洞见的可落地版）。发送方要让接收方 $j$ 复原任务相关变量 $Y$；$j$ 的边信息 $B_j=\mathcal M_j\ (\text{$j$ 可访问的记忆})+\text{廉价 ToM 估计}$。

\textbf{① 可达下界（WZ / Slepian--Wolf，解码端边信息）。} 接收方持有边信息 $B_j$（仅解码端）。\textbf{即使发送方不观测 $B_j$}，无损传 $Y$ 的可达速率为 $H(Y\mid B_j)$。这是地板。

\textbf{② 实用可达方案（预测残差编码）。} 发送方用 ToM 估计 $\hat B_j$ 形成预测 $\hat Y$，发残差 $Y-\hat Y$：
\[
R_{\text{SYNAPSE}}=H(Y\mid \hat B_j)=\underbrace{H(Y\mid B_j)}_{\text{WZ 地板}}+\underbrace{\Delta}_{\text{ToM 失配惩罚}},\qquad
\Delta\approx I(Y;B_j\mid \hat B_j)\ \ge 0.
\]
\begin{quote}\small
\textbf{诚实声明}：预测残差是一个\textbf{次优}的 WZ 方案（不是最优分箱）。我们不重证 WZ 逆定理；我们\textbf{应用并扩展}它到"可学习/增长边信息 + LLM-agent"。$\hat B_j=B_j$ 时 $\Delta=0$；$\hat B_j$ 与 $B_j$ 无关时 $\Delta=I(Y;B_j)$（边信息红利全丢）。\textbf{ToM 预测器的全部价值就是压低 $\Delta$。}
\end{quote}

\textbf{③ 收缩律（头牌现象，命题）。} 任务 $\sim\mathcal D$ 分段平稳，巩固算子使记忆朝 $\mathcal D$ 的充分统计收敛（$\mathrm{KL}(\mathcal M\,\|\,\mathrm{suff.stat}(\mathcal D))\downarrow$），则期望每任务协作速率
\[
\mathbb E_{T\sim\mathcal D}\big[H(Y\mid B_t)\big]\ \text{单调不增}\ \longrightarrow\ H(Y\mid B_\infty)\quad(\text{新信息熵地板}).
\]

\textbf{④ 非显然、可证伪的预测（让理论挣到饭钱）。}
\begin{quote}\small
收缩曲线\textbf{会停在一个可独立测量的地板} $H(Y_{\text{new}}\mid B_\infty)$。\textbf{怎么独立测？}给一个"全知接收方"喂入\textbf{全部巩固经验}，测它对新任务残差的熵——这就是预言的 plateau 高度。\textbf{曲线若真停在这条线上 $\to$ 理论被验证；若停不到 $\to$ 被证伪。}
\end{quote}

\textbf{⑤ 漂移回弹（应用，命题）。} $\mathcal D\to\mathcal D'$，巩固出的 $B_\infty$ 对 $\mathcal D'$ 不再充分，$H(Y'\mid B_\infty)>H(Y\mid B_\infty)$，\textbf{协作速率阶跃上升}；因 agent 会"多说几句补偿"，\textbf{回弹先于任务质量下降发生}。
\begin{quote}\small
\textbf{推论（可发表的应用）}：\textbf{协作速率是 MAS 分布漂移的一个免费、模型无关、提前的传感器}——不需要标签、不需要额外探针、不需要 ground truth。
\end{quote}

\subsection{与"缓存"的本质区别（全局最关键的一刀）}
这是评委和审稿人都会问的第一个问题，标准答案务必背下来：
\begin{quote}\small
\textbf{缓存吃"时间冗余"}：同一 query $\to$ 同一输出，精确前缀复用，省的是\textbf{重新生成的算力（compute）}。\\
\textbf{SYNAPSE 吃"语义冗余"}：跨 agent、面向\textbf{全新任务}，省的是\textbf{为对齐而必须交换的协作比特（coordination bits）}。\\
\textbf{两者正交可叠加。}在缓存之上，SYNAPSE 仍能在新任务上继续压低跨 agent 字节——因为新任务的答案从没被缓存过，但它的\textbf{结构}能被长大的记忆\textbf{预测}。
\end{quote}
一句话：\key{缓存让你"不必重算"，SYNAPSE 让你"不必重说"。}

\subsection{四个贡献（从 v4 的"4 并列"收敛到"1+1+1+1"）}
\begin{table}[h]\small\centering
\renewcommand{\arraystretch}{1.25}
\begin{tabularx}{\linewidth}{@{}p{0.8cm}p{1.8cm}Y p{1.7cm}@{}}
\toprule
 & \textbf{类型} & \textbf{内容} & \textbf{框架无关?} \\
\midrule
\textbf{C1} & 现象（头牌） & \textbf{协作速率收缩律 + 漂移回弹}：可测曲线、可独立测的地板、漂移阶跃 & 是（不信 WZ 也成立） \\
\textbf{C2} & 透镜（已被验证） & 形式化为\textbf{带增长边信息的语义 WZ}，给出可达方案 + 失配惩罚 $\Delta$；明确区分于 2604.09521（静态能力） & --- \\
\textbf{C3} & 保证（部署级） & \textbf{Verified Lossy Coordination}：校验和让有损语义信道\textbf{端到端正确}（代价 $O(\text{几十字节})$） & 是 \\
\textbf{C4} & 应用（so-what） & \textbf{协作速率作为模型无关的提前漂移/新颖性传感器} & 是 \\
\bottomrule
\end{tabularx}
\end{table}

\noindent 降级为机制/消融（保留但不当头牌）：言语行为决策（tell/ask/write/hold）、ToM 寻址、MDL 压缩层。

\section{赛题核心要求与"咬合"自检}

\begin{quote}\small
\textbf{强制自检（来自本项目历史教训 L9）}：收敛主创新后，必须逐条回扫赛题硬要求，确认每条都被覆盖、且\textbf{与主创新咬合}（而非孤立模块）。漏一条 = 验收红线失分。下表的 MUST 清单依据 v2/v4 整理，\textbf{以官方赛题书为最终准}。
\end{quote}

\subsection{硬要求 × 咬合点（MUST 清单）}
{\footnotesize
\renewcommand{\arraystretch}{1.22}
\begin{longtable}{@{}p{0.7cm}p{4.4cm}p{4.0cm}p{6.0cm}@{}}
\toprule
\textbf{\#} & \textbf{赛题硬要求} & \textbf{SYNAPSE 覆盖组件} & \textbf{与主创新的咬合} \\
\midrule
\endfirsthead
\toprule
\textbf{\#} & \textbf{赛题硬要求} & \textbf{SYNAPSE 覆盖组件} & \textbf{与主创新的咬合} \\
\midrule
\endhead
M1 & $\ge 3$ agent、$\ge 3$ 类角色、多步复杂任务 & Planner/Retriever/Executor·CodeAct/Summarizer（4 角色） & 角色间的对齐就是收缩律的测量对象 \\
M2 & 结构化通信 + 握手/能力发现/协议映射，禁纯文本透传 & 类型化消息 Schema + CNR 协商 & 残差句柄走结构化消息头 \\
M3 & 同时支持"纯文本"与"结构化协议"模式，可复现 A/B & 双模式开关 + 评测器 & A/B 正是收缩律 vs 基线的对照 \\
\textbf{M4} & \textbf{非文本中间状态传递（embedding/向量/隐状态），说明生成/传递/接收/使用} & \textbf{状态交换数据平面 = 预测残差编码}（见 §4.2） & \textcolor{AcademicRed}{\textbf{★ 核心咬合：非文本载体就是"惊讶残差"，不是"发更大的向量"}} \\
M5 & 共享记忆单元（ID/来源/时间/主题/摘要等元数据） & 统一记忆单元 schema & 记忆 = 解码端边信息 $B_j$ 的物理载体 \\
M6 & 关键词/标签/语义检索，跨任务复用 & 倒排 + 向量混合检索 & 检索命中越多 $\to B_j$ 越大 $\to$ 速率越低（收缩的微观机制） \\
M7 & $\ge 2$ 组关联连续任务验证 & 任务族 G1/G2（见 §5.2） & 关联性 = 任务间共享结构 = 收缩的前提 \\
M8 & 全套统计（消息数/token/非文本字节/耗时/命中率/提升） & 评测模块 + eBPF 探针 & 收缩曲线 = 这些统计随轮次的轨迹 \\
M9 & 五大模块齐备、稳定 $\ge 10$ 轮 & 五模块架构（见 §4） & $\ge 10$ 轮正好画出收缩曲线 \\
M10 & openEuler 可编译运行；交付源码/文档/视频 & 全栈容器化 & 现场 demo = 收缩-回弹仪表盘 \\
\bottomrule
\end{longtable}}

\begin{quote}\small
\textcolor{AcademicRed}{\textbf{★ M4 是最容易被"聚焦主线"挤出视野、却最关键的一条}}（L9 的原始教训就是这条被漏）。SYNAPSE 的设计让 M4 与主创新\textbf{天然咬合}：非文本中间状态传递的机制\textbf{就是}预测残差编码，\textbf{不是}额外贴一个"发 embedding"的模块。这既补齐验收红线，又是一个差异化加分点。
\end{quote}

\subsection{评分细则打法（100 分）}
\begin{table}[h]\small\centering
\renewcommand{\arraystretch}{1.25}
\begin{tabularx}{\linewidth}{@{}p{2.6cm}p{1.0cm}Y@{}}
\toprule
\textbf{评分项} & \textbf{分值} & \textbf{打法} \\
\midrule
通信效率 & 25 & 预测残差 + 引用传递，token 与字节双降；\textbf{eBPF 实测}坐实 \\
状态传递创新 & 20 & 非文本 = "发更小的惊讶" + \textbf{Verified Lossy 校验}（有损还保证正确，评委能懂其工程价值） \\
记忆复用 & 20 & 记忆即边信息；收缩曲线\textbf{就是}记忆复用增益的可视化证据 \\
系统完整性 & 20 & 五模块、4 角色、CodeAct 沙箱、$\ge 10$ 轮、openEuler 跑通 \\
实验验证 & 15 & 双模式 A/B + 关联任务 + 因果消融 + \textbf{诚实 crossover}（写明"何时不该用"反增可信度） \\
\bottomrule
\end{tabularx}
\end{table}

\section{系统架构}

\subsection{五大模块（一张图看懂）}
可编辑版见同目录 \texttt{SYNAPSE\_architecture.drawio}。逻辑结构如图~\ref{fig:arch}：自上而下的五层 + openEuler 底座，右侧虚线是\textbf{核心反馈环}——巩固使记忆长大、ToM 把 $\hat B_j$ 估得更准、残差更小，于是\textbf{协作速率沿曲线收缩}。

\begin{figure}[h]\centering
\begin{tikzpicture}[node distance=4mm]
\node[band, fill=blue!7] (l1) {\textbf{① 多 Agent 运行时}\quad（开源同族模型 / vLLM，可取隐状态）\\[1pt]
Planner（规划） \;·\; Retriever（检索） \;·\; Executor·CodeAct（沙箱执行） \;·\; Summarizer（汇总）};
\node[band, fill=orange!8, below=of l1] (l2) {\textbf{② 协议解析与调度}\\[1pt]
类型化消息 Schema \{action,params,handles,capability\} \;·\; 握手/能力发现 CNR \;·\; \textbf{言语行为选择器 tell/ask/write/hold}};
\node[band, fill=yellow!20, below=of l2] (l3) {\textbf{③ 状态交换·数据平面（核心 = 非文本状态传递）}\\[1pt]
内容寻址库 CAS \;·\; 接收方信念 $B_j$（记忆+ToM 估计）\;·\; \textbf{预测残差编码（只发 $H(Y\mid B_j)$ 的"惊讶"）}\;·\; \textbf{校验和·Verified Lossy（hash 不符$\to$回退 ask(full)）}\;·\; 三级表示：隐状态/句向量/文本 diff};
\node[band, fill=green!9, below=of l3] (l4) {\textbf{④ 共享记忆与认知}\\[1pt]
记忆单元+元数据（ID/来源/时间/主题/摘要）\;·\; 混合检索（关键词/标签/语义）\;·\; \textbf{ToM 预测器}（估计 $B_j$，重构成功率自监督）\;·\; \textbf{巩固器（睡眠/回放）}};
\node[band, fill=black!6, below=of l4] (l5) {\textbf{⑤ 评测与度量}\\[1pt]
双模式 A/B（text vs synapse）\;·\; eBPF 字节/时延/质量/FLOPs \;·\; \textcolor{NavyBlue}{\textbf{收缩-回弹仪表盘：$R\downarrow$ 逼近地板 $\to$ 注入漂移 $\to R\uparrow$ 回弹}}};
\node[band, fill=violet!8, below=of l5] (l6) {\textbf{底座 · openEuler 24.03-LTS-SP3}\\[1pt]
vLLM \;|\; 向量库 (Milvus/FAISS) \;|\; IPC·共享内存/mmap·Socket·gRPC \;|\; WASM/容器沙箱 \;|\; eBPF 探针};
\draw[flow] (l1)--(l2); \draw[flow] (l2)--(l3); \draw[flow] (l3)--(l4); \draw[flow] (l4)--(l5); \draw[flow] (l5)--(l6);
\draw[-{Stealth[length=2.6mm]}, dashed, line width=0.9pt, NavyBlue]
  (l4.east) to[out=0,in=0,looseness=2.2] node[right, font=\footnotesize, text=NavyBlue]{经验$\uparrow\,\to$收缩} (l3.east);
\end{tikzpicture}
\caption{SYNAPSE 系统架构：五层 + 底座，右侧虚线为"经验增长$\to$速率收缩"核心反馈环。}
\label{fig:arch}
\end{figure}

\subsection{非文本状态传递：生成 / 传递 / 接收 / 使用（赛题 M4 点名要求）}
以"Retriever $\to$ Summarizer 传递检索得到的长证据"为例：
\begin{itemize}
\item \key{生成}：Retriever 取末层隐状态摘要（或对片段编码出句向量）作为任务相关表示 $Y$；用 ToM 估计 Summarizer 已能预测的部分 $\hat Y(\hat B_j)$；算残差 $z=Y-\hat Y$。残差是\textbf{低熵的}（只剩惊讶）$\to$ 可激进量化/稀疏化 $\to$ \textbf{字节真小}（区别于 CIPHER/KV 发全量向量）。
\item \key{传递}：消息 = 类型化头 \{action,params,capability\} + \textbf{残差句柄/短码} + $\mathrm{hash}(Y)$；CASD 只传 $j$ 缺的差集。
\item \key{接收}：Summarizer 按 CNR 协商解码（同族$\to$隐状态零拷贝注入；异构$\to$取该句柄的句向量/文本视图）；用 $B_j$ 预测 + 残差\textbf{重构} $\hat Y$；\textbf{校验 $\mathrm{hash}(\hat Y)=\mathrm{hash}(Y)$？不符 $\to$ ask(full)}。
\item \key{使用}：作为软前缀/检索条件继续生成，省掉"内部态$\to$文本$\to$内部态"往返。
\item \key{诚实回退}：若残差不比文本短、或对方异构不可解码，言语行为选择器自动选文本/hold，由 eBPF 记录该 crossover。
\end{itemize}
\begin{quote}\small
这一节把 v1 审计 A5（"非文本常更大"）和赛题 M4 一次解决：\textbf{非文本在这里不是"发更大的向量"，而是"发更小的惊讶"}——这是字节级真省的关键。
\end{quote}

\subsection{L0/L1 分层（一套系统，两层产出，demo 去风险）}
\begin{table}[h]\small\centering
\renewcommand{\arraystretch}{1.25}
\begin{tabularx}{\linewidth}{@{}p{1.5cm}Y p{2.0cm}p{3.0cm}@{}}
\toprule
\textbf{层} & \textbf{内容} & \textbf{服务对象} & \textbf{稳定性} \\
\midrule
\textbf{L0 零阶} & 接收方信念 $\approx$ 记忆视图；残差 = 记忆差集（+连续表示锚点残差）；言语行为 = \textbf{规则打分}（差集大$\to$tell；自身缺$\to$ask；高价值中间态$\to$write；残差$\approx0\to$hold）；\textbf{校验和} & \textbf{竞赛现场 demo} & \textbf{全确定性，不依赖会抖的学习件，不会崩} \\
\textbf{L1 学习层} & 轻量 ToM 预测器（小 MLP/检索式，重构成功率自监督）；EFE 打分器（离线监督/bandit，非在线 RL）；巩固器；MDL 码本；收缩律拟合 & \textbf{顶会论文} & \textbf{离线日志展示，现场不实时跑} \\
\bottomrule
\end{tabularx}
\end{table}

\section{实验设计与预测花销}

\subsection{实验清单（每个实验都对应一个待拆的攻击）}
{\footnotesize
\renewcommand{\arraystretch}{1.22}
\begin{longtable}{@{}p{1.6cm}p{2.6cm}p{5.2cm}p{5.6cm}@{}}
\toprule
\textbf{编号} & \textbf{目的} & \textbf{设计} & \textbf{预期结果 / 拆哪个攻击} \\
\midrule
\endfirsthead
\toprule
\textbf{编号} & \textbf{目的} & \textbf{设计} & \textbf{预期结果 / 拆哪个攻击} \\
\midrule
\endhead
\textbf{E1（头牌）} & 收缩律 + 地板 & $\ge2$ 组关联连续任务、$\ge10$ 轮，画 $R$（每任务跨 agent 字节）vs 累计经验；拟合形态/拐点/地板；\textbf{全知接收方 oracle 独立测地板}；消融巩固/ToM $\to$ 收缩消失 & 单调收缩、停在 oracle 地板（\textbf{可证伪}）；用\textbf{新任务}、且与缓存共存 $\to$ 证明是协作速率收缩而非答案缓存（拆"就是缓存"） \\
\textbf{E2} & 与缓存正交 & 缓存 on/off $\times$ SYNAPSE on/off 四象限 & 增益叠加（拆"缓存已 90\%"） \\
\textbf{E3} & 失配惩罚 $\Delta$ & 人为劣化 ToM $\to$ 看 $\Delta$、字节$\uparrow$、校验回退率$\uparrow$ & $\Delta$ 随 ToM 质量单调（验 §2.2 ②） \\
\textbf{E4} & 正确性 & 注入 ToM 错误，校验和保证端到端正确 & 任务正确性不被有损残差破坏（拆"静默损坏"） \\
\textbf{E5} & 成本 crossover & 消息长度 $\times$ 同/异构 $\times$ 预测器开销的净收益相图（bytes+FLOPs+latency） & 给出"何时净赢/净亏"；选择器自动选文本（拆"字节换 FLOPs"） \\
\textbf{E6} & 漂移回弹（应用） & 注入分布切换，观察 $R$ 回弹再收缩；与"质量下降时刻"对齐 & 回弹\textbf{先于}质量下降 $\to$ 漂移传感器（拆"平稳性假设"，把局限变卖点） \\
\textbf{E7} & 言语行为消融 & 去 ask/hold/认知项 & 冗余轮次与字节上升（验认知项真有用） \\
\textbf{E8} & 表示三级 & 隐状态/句向量/文本 diff 的省幅与适用域 & 同族省最多，异构文本-diff 仍省（拆"异构即崩"） \\
\bottomrule
\end{longtable}}

\subsection{任务族设计（收缩的前提是"有共享结构"）}
收缩幅度取决于任务间共享结构（这是诚实风险 R1）。因此\textbf{任务族必须刻意构造共享结构，并把"结构量 $\to$ 收缩幅度"测成一条结论曲线}（横轴用任务间互信息/主题重叠度的代理量）：
\begin{itemize}
\item \key{G1（主题深挖）}：对主题 X 做多跳研究并产出报告（plan$\to$retrieve$\to$CodeAct 计算$\to$summarize），跑若干轮。
\item \key{G2（关联演进）}：X 的更新/对比/再分析，\textbf{复用 G1 的证据句柄 + 计划过程}。度量复用带来的"少检索、少计算、少 deliberation"。
\item \key{漂移臂}：在 G2 中途切到与 X 弱相关的主题 X$'$，触发 E6 回弹。
\item \key{对照族（负例）}：故意选\textbf{无共享结构}的任务流，预期收缩幅度 $\approx 0$ —— 用它证明"收缩来自共享结构"而非系统假象。
\end{itemize}

\subsection{度量（模式无关主指标 + 分量诚实并报）}
\begin{itemize}
\item \key{主指标}：端到端\textbf{时延}与\textbf{能量 (J)}（模式无关，避免 token vs 字节"苹果比橘子"）。
\item \key{分量}：$\mathrm{bytes}_{\text{wire}}$（\textbf{eBPF 实测}）$\oplus$ FLOPs(ToM+编解码)（计数）$\oplus$ latency $\oplus$ energy（功率模型\textbf{估算}，标注 estimated）。
\item \key{质量护栏}：LLM-judge / ground truth，证明"省了但不掉质量"。\textbf{校验和保证正确性，质量门是双保险}。
\end{itemize}

\subsection{预测花销}
\begin{quote}\small
\textbf{算力预算 source: assumed}（未与你确认硬件）。下表按"单台 24--40GB GPU + 开源同族模型 (Qwen2.5-7B/14B-Instruct，vLLM)"估算；若你有更强机器可线性缩短。\textbf{此项进首个审批决策包，请你校准。}
\end{quote}
\begin{table}[h]\small\centering
\renewcommand{\arraystretch}{1.22}
\begin{tabularx}{\linewidth}{@{}p{3.0cm}p{4.2cm}Y@{}}
\toprule
\textbf{项目} & \textbf{估算} & \textbf{说明} \\
\midrule
模型 & Qwen2.5-7B + 14B-Instruct（同族，可取隐状态） & 同族是隐状态零拷贝的前提；异构臂用句向量/文本 \\
任务-运行规模 & $\sim$2 族 $\times\sim$80 任务 $\times\sim$10 轮 $\times$ \{text, synapse\} $\times\sim$6 消融 $\times$ 3 seed & $\approx$ 1{,}000--1{,}500 次任务-运行 \\
token 体量 & $\approx$ 50--80 M tokens（推理，全本地） & 多 agent $\times$ 多轮，$\sim$50K tokens/任务-运行 \\
GPU-小时 & \textbf{$\approx$ 30--60 GPU·h}（含重跑/seed） & 7B 在 4090 约 2K tok/s；一遍 $\sim$8--10h，$\times$ 重跑 \\
ToM 预测器训练 & 可忽略（小 MLP/检索式，离线，CPU/分钟级） & 监督信号 = 重构成功率，无需大训练 \\
质量评审 API & \textbf{$\approx$ \$20--80}（LLM-judge，可选） & 也可用同族模型自评，进一步省钱 \\
eBPF/系统适配 & 工程工时，非算力 & 提前在 openEuler 24.03 验证工具链 \\
总钱预算 & \textbf{$\approx$ \$20--80 + 自有 GPU 电费}；租云约 \$30--100 & 整个项目级，非每次实验 \\
总墙钟 & \textbf{$\approx$ 8--12 周}（M1--M5，单人/小队，含工程） & 见 §6.1 里程碑 \\
\bottomrule
\end{tabularx}
\end{table}

\noindent\textbf{结论}：这是一个"小钱 + 一台 GPU + 两三个月"就能跑出全套可证伪曲线的题——这正是它对竞赛+投稿双周期友好的关键（对照 v1 的在线多智能体 RL：那是算力黑洞，已在 v2 审计判死）。

\section{执行计划}

\subsection{里程碑（M1--M5）}
\begin{enumerate}
\item \key{M1（2 周）· 骨架}：运行时 + 类型化协议 + 双模式 + 评测器骨架；先跑通 text 基线。\textbf{门禁}：text 模式稳定 $\ge10$ 轮。
\item \key{M2（2--3 周）· 数据平面}：CAS + mmap 隐状态 + 向量库 + CNR 协商 + \textbf{预测残差编码 + 校验和}（L0 全确定性）。\textbf{门禁}：synapse 模式字节 $<$ text，且校验回退率可测；出 E4/E2 初值。
\item \key{M3（2--3 周）· 记忆与认知}：共享记忆 + 混合检索 + 巩固器 + \textbf{ToM 预测器（L1）}。\textbf{门禁}：跑出 E1 收缩曲线第一版 + oracle 地板。
\item \key{M4（2 周）· 实验闭环}：CodeAct 沙箱 + eBPF 度量 + G1/G2/漂移臂 + 全消融（E3/E5/E6/E7/E8）。\textbf{门禁}：E1 收缩 + E6 回弹 + E2 正交三张核心图成立。
\item \key{M5（2 周）· 交付}：源码/设计文档/部署文档/实验报告/演示视频；openEuler 复现；收缩-回弹仪表盘 demo 定版。（投稿向：补 L1 分析 + 论文初稿。）
\end{enumerate}
\begin{quote}\small
\textbf{关键路径排序原则}：先把 E1（收缩+地板）和 E2（缓存正交）设计到可证伪、可复现，\textbf{再大规模铺工程}——这两个实验同时决定"论文成不成"和"竞赛哇点立不立"，是命门。
\end{quote}

\subsection{风险登记册（不粉饰）}
\begin{table}[h]\small\centering
\renewcommand{\arraystretch}{1.22}
\begin{tabularx}{\linewidth}{@{}p{4.0cm}p{1.0cm}Y@{}}
\toprule
\textbf{风险} & \textbf{等级} & \textbf{解法} \\
\midrule
R1 收缩幅度小（共享结构少时 $H(Y\mid B)\approx H(Y)$） & 高 & 刻意构造有共享结构的任务族；把"结构量$\to$收缩幅度"\textbf{测成结论}；负例族证因果 \\
R2 ToM 无信念 ground truth & 中 & 不预测"信念"，预测\textbf{接收方实际重构成功率}（校验和成败做自监督） \\
R3 WZ "只是类比" & 中 & 给实用可达方案 + 两个可测命题（$\Delta$、地板）+ 明说"应用并扩展，不重证逆定理"；2604.09521 反证此透镜被学界认可 \\
R4 能量度量繁琐 & 低 & 主报 bytes(eBPF)+latency 硬数据，energy 用功率模型估算并标 estimated \\
异构即崩 & 中 & 表示三级（隐状态/句向量/文本 diff）优雅降级；WZ 框架表示无关 \\
demo 现场崩 & 中 & L0 全确定性现场跑，L1 离线出图，不让学习件上现场 \\
openEuler 适配 & 中 & 全栈容器化；M1 前先验证工具链 \\
\bottomrule
\end{tabularx}
\end{table}

\subsection{交付物清单}
\begin{itemize}
\item \key{竞赛}：源码 + 设计文档 + 部署文档（openEuler 24.03 容器化）+ 实验报告（全套赛题统计 + crossover）+ 演示视频（收缩-回弹仪表盘）。
\item \key{论文}：本方案 $\to$ idea-card $\to$ 实验协议 $\to$ 分析报告 $\to$ 顶会初稿（题目候选见附录 C）。
\item \key{一套系统两层产出}：先交 L0 拿竞赛分与可复现实验 $\to$ 加 L1 学习层与定律分析成稿投顶会。
\end{itemize}

\section{定位与新颖性（撞车窗检索结论）}
检索工具：semantic-scholar + web-search-prime（2026-06-17 执行）。下列论文\textbf{已逐一核到存在}并确认其真实主张（除标注外）。

{\scriptsize
\renewcommand{\arraystretch}{1.25}
\begin{tabularx}{\linewidth}{@{}p{1.4cm}YYYYY Y@{}}
\toprule
\textbf{维度} & \textbf{缓存族} & \textbf{潜空间通信} & \textbf{2604.09521 语义率失真} & \textbf{C2C (2510.19995)} & \textbf{Scaling Teams/Time} & \textbf{SYNAPSE} \\
\midrule
省什么 & 重复再生成 & 发全量潜表示 & （理论刻画字母表） & 建模协作成本 & 性能 vs 团队/时间 & \textbf{跨 agent 协作速率} \\
边信息 & 精确前缀 & 无 & 静态能力异质 & 共享心智(描述) & --- & \textbf{共享记忆+ToM，随经验增长} \\
时间动力学 & 命中率随 warmth & 无 & \textbf{无} & 无 & 性能随终身学习 & \textbf{速率收缩 + 漂移回弹（新）} \\
正确性 & 精确/近似 & 会静默损坏 & --- & --- & --- & \textbf{有损残差 + 校验$\to$正确} \\
系统/度量 & 应用级 & token & \textbf{纯理论} & 时延 & 性能 & \textbf{真系统 + eBPF} \\
与缓存 & --- & 不正交 & --- & --- & --- & \textbf{正交可叠加} \\
\bottomrule
\end{tabularx}}

\medskip
\noindent\textbf{与最近邻 2604.09521 的精确区分}（这是 SYNAPSE 的护城河）：它做\textbf{静态、纯理论、能力异质}的率失真刻画；我们做\textbf{动态（边信息随经验增长）、真系统、字节级实证}的\textbf{时间收缩 + 漂移回弹 + 校验有损 + ToM 寻址}。它把透镜立了起来——正好替我们去掉"这个框架合不合法"的风险——但\textbf{金矿（时间动力学 + 系统）它一寸没挖}。

\section*{附录 A · 诚实声明（assumed 项 / 待核验项）}
\addcontentsline{toc}{section}{附录 A · 诚实声明}
\begin{itemize}
\item \key{算力预算 source: assumed}：§5.4 按单台 24--40GB GPU + Qwen2.5-7B/14B 估算，未与你确认硬件 $\to$ 进首个审批决策包。
\item \key{赛题 MUST 清单 source: reconstructed}：§3.1 依据 v2/v4 整理，\textbf{以官方赛题书为最终准}；提交前请逐条对官方原文。
\item \key{待核验 novelty（投稿前必补一次定向检索）}：
  \begin{itemize}
  \item "ToM 用于\textbf{最小消息寻址}"是本轮核验\textbf{最弱}的一条（ToM-MAS 都在追踪信念，未见用于压消息字节，但未穷尽）$\to$ 投稿前检索 audience design / pragmatic compression / RSA $\times$ agents。
  \item \textbf{arXiv:2511.04235}（v4 称"空间预测编码 + HRL"）本环境工具受限\textbf{未独立核到域}，请自查其域与是否用 RL，再决定引用措辞。
  \end{itemize}
\item \key{方法定位诚实}：预测残差是\textbf{次优} WZ 方案；我们不声称达到 WZ 最优，也不声称新逆定理。
\end{itemize}

\section*{附录 B · 关键参考（★ = 本轮已核到存在并确认主张）}
\addcontentsline{toc}{section}{附录 B · 关键参考}
\begin{itemize}\small
\item ★ Semantic Rate-Distortion for Bounded Multi-Agent Communication. arXiv:2604.09521（最近邻：静态能力角度，须正面区分）
\item ★ Communication to Completion (C2C). arXiv:2510.19995（建模协作成本 + Alignment Factor）
\item ★ Agentic Plan Caching: Test-Time Memory for Fast and Cost-Efficient LLM Agents. arXiv:2506.14852（计划缓存，须正交化区分）
\item ★ Scaling Teams or Scaling Time? Memory-Enabled Lifelong Learning in LLM MAS. arXiv:2604.03295（性能-时间权衡）
\item ★ Collaborative Memory: Multi-User Memory Sharing with Dynamic Access Control. arXiv:2505.18279（ICML 2025）
\item Wyner--Ziv (1976) 解码端边信息有损信源编码；Slepian--Wolf (1973) 分布式无损编码（经典基底）
\item CIPHER. arXiv:2310.06272 \;|\; KVComm. arXiv:2510.03346 \;|\; LatentMAS+OBF（须核 ID）
\item Theory of Mind for Multi-Agent Collaboration via LLMs. EMNLP 2023 \;|\; Mutual ToM. arXiv:2409.08811
\item Shared Spatial Memory Through Predictive Coding. arXiv:2511.04235（域待核）
\item AgentPrune / Cut the Crap. arXiv:2410.02506（ICLR'25，空间轴 scaling 对照）
\end{itemize}

\section*{附录 C · 题目候选}
\addcontentsline{toc}{section}{附录 C · 题目候选}
\begin{itemize}\small
\item \emph{Coordination as Compression: A Contraction Law for Multi-Agent LLM Communication with Growing Memory}
\item \emph{Speak Less as You Remember More: Verified Wyner--Ziv Coordination and a Drift Sensor for LLM Agents}
\item \emph{The Communication Contraction Law: How Shared Memory Shrinks---and Drift Rebounds---Multi-Agent LLM Bandwidth}
\end{itemize}

\vspace{4mm}
\noindent\rule{\linewidth}{0.4pt}\\[1mm]
{\footnotesize\itshape 收敛要点：头牌从"语义 WZ"换成"协作速率时间律 + 漂移传感器"，WZ 降为可引用透镜（避开 2604.09521 的正面撞车，占据其未挖的动力学+系统无人区）；新增 Verified Lossy 命名贡献、地板可证伪预测、ToM 重构成功率自监督；L0/L1 分层保证竞赛稳跑、论文出曲线。}

\end{document}
